%! The Complete Solution to Ax = b — Unit 3, Lesson 3.3 — Warm-Up (Answer Key)
\documentclass[10pt]{article}
\usepackage{linalg-article}
\usepackage{linalg-key}   % pulls in -boxes; do NOT also load -boxes
\newcommand{\TallMath}[1]{$\displaystyle #1\rule[-1.4em]{0pt}{3.2em}$}
\newcommand{\xx}{\mathbf{x}}
\newcommand{\bb}{\mathbf{b}}
\newcommand{\svec}{\mathbf{s}}
\newcommand{\zero}{\mathbf{0}}

\begin{document}

\pageheader{Unit 3, Lesson 3.3}{Warm-Up --- Answer Key}
\namedateperiod

\vspace{0.12in}

\begin{notesbox}{Getting Ready: Skills We'll Use Today}
\small These skills from Units 1--2 and Lesson 3.2 power today's work. Answer quickly.

\vspace{4pt}
\textbf{1. Find a nullspace (Lesson 3.2).} Reduce $A=\begin{bmatrix} 1 & 2 \\ 2 & 4 \end{bmatrix}$ and
give the special solution of $A\xx=\zero$:
\[
  \begin{bmatrix} 1 & 2 \\ 2 & 4 \end{bmatrix}
  \;\longrightarrow\;
  \begin{bmatrix} 1 & 2 \\[2pt] {\color{keyred}\mathbf{0}} & {\color{keyred}\mathbf{0}} \end{bmatrix},
  \qquad
  \svec=\begin{bmatrix}\rule{0pt}{1.1em}{\color{keyred}\mathbf{-2}}\\[2pt]{\color{keyred}\mathbf{1}}\end{bmatrix}.
\]

\vspace{4pt}
\textbf{2. Solve $A\xx=\bb$ by elimination (Unit 2).} Solve
$\begin{bmatrix} 1 & 1 \\ 1 & 2 \end{bmatrix}\begin{bmatrix} x_1 \\ x_2 \end{bmatrix}=\begin{bmatrix} 3 \\ 5 \end{bmatrix}$:
\[
  x_1=\ans{$1$}, \qquad x_2=\ans{$2$}.
\]

\vspace{4pt}
\textbf{3. Can we reach the target? (Lesson 1.3)} The columns of $A=\begin{bmatrix} 1 & 3 \\ 2 & 6 \end{bmatrix}$
both lie on the line $y=2x$. Is the target $\bb=\begin{bsmallmatrix} 2\\4 \end{bsmallmatrix}$ a combination
of these columns --- that is, does $A\xx=\bb$ have a solution? Circle one and give a one-line reason.
\quad \ans{yes} / no
\ansline{Yes --- $\bb=\begin{bsmallmatrix} 2\\4 \end{bsmallmatrix}$ is on the line $y=2x$ (e.g.\ $\bb=2\begin{bsmallmatrix} 1\\2 \end{bsmallmatrix}$), so it \emph{is} a combination of the columns: $\bb\in C(A)$.}
\end{notesbox}

\begin{teachernote}
Item~1 supplies the \emph{nullspace part} we add today ($\svec=\begin{bsmallmatrix} -2\\1 \end{bsmallmatrix}$).
Item~2 is the elimination that produces a \emph{particular solution} $\xx_p$. Item~3 is the
\emph{solvability} idea in disguise: $A\xx=\bb$ has a solution exactly when $\bb$ lies in $C(A)$ (here the
line $y=2x$). Today's lesson snaps these three together: $\xx=\xx_p+\xx_n$, provided $\bb$ is reachable.
\end{teachernote}

\end{document}
